\exnumber{02}
\extitle{UX II - Ex 3 - Affinity diagram \& Empathy map}
\exfiles{ex02/zip archive, named Name\_FirstName\_DeliverableName\_Date\_VersionNumber}
\exauthorize{Those named in the exercise}
\makeheaderfiles

\textbf{Instructions:}

Open your written scripts and extract the relevant information.

1. Use paper or digital sticky notes (Post-its or Miro.com) to isolate the pieces of information your interviewees shared with you.
1 idea, quote, or piece of information = 1 sticky note.
2. When all your scripts are written/typed down on sticky notes, gather them in clusters based on themes, opinions, habits, or anything they could have in common.
3. Label the clusters.
4. Draw conclusions
5. Fill an empathy map ONLY with the content from your affinity diagram.

\textbf{Deliverables:}

\begin{itemize}
\item An Affinity Diagram (pictures of the Affinity Diagram on paper or on digital tools like Miro)
\item An Empathy Map (pictures of the Empathy Map on paper or on digital tools like Miro)
\item Upload the documents on Github, in a zip folder named "ProjectTitle\_Name\_FirstName", with all project deliverables as follows: "Name\_FirstName\_DeliverableName\_Date\_VersionNumber".  For example, the first deliverable can be named as "Doe\_John\_MidFiPrototype\_05242024\_V1".
\item Add a title within all written documents.
\item Pay attention to the global aspect of the deliverables. They must be clear, simple and easy to read. You can get inspiration on canva.com but don't overload your design with too much details!
\end{itemize}

\textbf{Resources:}

\begin{itemize}
\item What is an Affinity Diagram and how to use it \url{https://miro.com/blog/create-affinity-diagrams/}
\item Affinity Diagram and Relation Examples \url{https://www.youtube.com/watch?v=VngWHIE4k9s}
\item Updates Empathy Map Canvas by Dave Gray \url{https://medium.com/the-xplane-collection/updated-empathy-map-canvas-46df22df3c8a}
\item Empathy Map by Krisztina Szerovay \url{https://uxknowledgebase.com/empathy-map-ca037e7686b6}
\end{itemize}

\textbf{Quote:}

Remember, you are not the user, so your opinions are not absolute truth.

